\begin{styletable}

Отображение (элемент)
    & $x \mapsto f(x)$
    & \verb|\mapsto|
\\

Отображение (множество)
    & $f \colon X \to Y$
    & \verb|f \colon X \to Y|
\\

Образ подмножества
    & $f(A)$
    & \verb|f(A)|
\\

Прообраз подмножества
    & $f^{-1}(B)$
    & \verb|f^{-1}(B)|
\\

Прообраз элемента
    & $f^{-1}(\set{y})$
    & \verb|f^{-1}(\set{y})|
\\

Обратная функция
    & $f^{-1} \colon Y_f \to X$, $f^{-1}(y)$
    & \verb|f^{-1} \colon Y_f \to X|, \verb|f^{-1}(y)|
\\

Композиция
    & $g \circ f$
    & \verb|\circ|
\\

Тождественное равенство
    & $f_1 \equiv f_2$
    & \verb|\equiv|
\\

\hline
\end{styletable}

\stylenote{
    Запись \texttt{f\^{}\{-1\}(y)} в книге используется только для значения обратной функции.
    Для полного прообраза элемента используется запись $f^{-1}(\set{y})$.
    Таким образом, если $f$ биективна, то $f^{-1}(y)$~--- это элемент, а $f^{-1}(\set{y})$~--- одноэлементное множество $\set{f^{-1}(y)}$.
}